% !TEX program = pdflatex
% !TEX encoding = UTF-8 Unicode

% Plantilla, basada en la clase `scrbook` del paquete KOMA-script,  para la elaboración de un TFG siguiendo las directrices del la comisión del Grado en Matemáticas de la Universidad de Granada.

% Francisco Torralbo Torralbo

\documentclass{ugrTFG}
% !TeX root = tfg.tex
% !TeX encoding = utf8
%
%*******************************************************
% Preámbulo propio: paquetes y macros añadidos a la clase ugrTFG
% (se carga justo después de \documentclass; no modificar ugrTFG.cls)
%*******************************************************

% -------------------------------------------------------------------
% Versión borrador / final
% -------------------------------------------------------------------
% Con \borradortrue se muestran las marcas [PENDIENTE] y [REVISAR];
% con \borradorfalse desaparecen en la versión final.
\newif\ifborrador
\borradortrue

\ifborrador
  \newcommand{\pendiente}[1]{\textcolor{red}{\textbf{[PENDIENTE: #1]}}}
  \newcommand{\revisar}[1]{\textcolor{orange}{\textbf{[REVISAR: #1]}}}
\else
  \newcommand{\pendiente}[1]{}
  \newcommand{\revisar}[1]{}
\fi

% -------------------------------------------------------------------
% Paquetes
% -------------------------------------------------------------------
\usepackage{siunitx}
\sisetup{output-decimal-marker={,}, group-separator={.}, group-minimum-digits=4}
\usepackage{subcaption}
\usepackage{multirow}
\usepackage{tabularx}
\usepackage{longtable}
\usepackage{array}
\usepackage{pgfgantt}

% cleveref siempre después de hyperref (que carga la clase)
\usepackage[spanish, noabbrev]{cleveref}
\crefname{teorema}{teorema}{teoremas}
\Crefname{teorema}{Teorema}{Teoremas}
\crefname{proposicion}{proposición}{proposiciones}
\Crefname{proposicion}{Proposición}{Proposiciones}
\crefname{lema}{lema}{lemas}
\Crefname{lema}{Lema}{Lemas}
\crefname{corolario}{corolario}{corolarios}
\Crefname{corolario}{Corolario}{Corolarios}
\crefname{definicion}{definición}{definiciones}
\Crefname{definicion}{Definición}{Definiciones}
\crefname{observacion}{observación}{observaciones}
\Crefname{observacion}{Observación}{Observaciones}
\crefname{ejemplo}{ejemplo}{ejemplos}
\Crefname{ejemplo}{Ejemplo}{Ejemplos}
\crefname{part}{parte}{partes}
\Crefname{part}{Parte}{Partes}

% -------------------------------------------------------------------
% Entornos
% -------------------------------------------------------------------
% La clase no define "demostracion": se construye sobre proof de amsthm,
% que con babel en español ya imprime «Demostración».
\newenvironment{demostracion}{\begin{proof}}{\end{proof}}

% -------------------------------------------------------------------
% Notación recurrente (definida una sola vez)
% -------------------------------------------------------------------
% Conjuntos y probabilidad
\newcommand{\R}{\mathbb{R}}
\newcommand{\N}{\mathbb{N}}
\DeclareMathOperator{\E}{\mathbb{E}}
\newcommand{\Prob}{\mathbb{P}}
\newcommand{\Ind}{\mathbb{1}}
\DeclareMathOperator*{\argmax}{arg\,max}
\DeclareMathOperator*{\argmin}{arg\,min}
\DeclareMathOperator{\softmax}{softmax}
\DeclareMathOperator{\Var}{Var}

% Distribuciones y conjuntos ID / OOD
\newcommand{\Pin}{P_{\mathrm{in}}}
\newcommand{\Pout}{P_{\mathrm{out}}}
\newcommand{\Din}{\mathcal{D}_{\mathrm{in}}}
\newcommand{\Dout}{\mathcal{D}_{\mathrm{out}}}
\newcommand{\Xin}{X_{\mathrm{in}}}
\newcommand{\Xout}{X_{\mathrm{out}}}

% Métricas (válidas en texto y en modo matemático)
\newcommand{\AUROC}{\ensuremath{\mathrm{AUROC}}\xspace}
\newcommand{\AUPRin}{\ensuremath{\mathrm{AUPR\text{-}In}}\xspace}
\newcommand{\AUPRout}{\ensuremath{\mathrm{AUPR\text{-}Out}}\xspace}
\newcommand{\FPR}{\ensuremath{\mathrm{FPR}}\xspace}
\newcommand{\TPR}{\ensuremath{\mathrm{TPR}}\xspace}
\newcommand{\CCR}{\ensuremath{\mathrm{CCR}}\xspace}
\newcommand{\FPRn}[1]{\ensuremath{\mathrm{FPR@}#1}\xspace}
\newcommand{\CCRn}[1]{\ensuremath{\mathrm{CCR@FPR}#1}\xspace}

% Nombres de métodos y conjuntos (formato consistente en el texto)
\newcommand{\stoodx}{STOOD-X\xspace}
\newcommand{\nctcrc}{NCT-CRC-HE-100K-NONORM\xspace}
\newcommand{\clase}[1]{\texttt{#1}}

\endinput
 % Paquetes y macros propios (no modificar ugrTFG.cls)

% VERSIÓN ELECTRÓNICA PARA TABLETA
% Activando la opción "tablet" generaremos un pdf adaptado para leerlo en tabletas de 9 pulgadas:
% \documentclass[tablet]{ugrTFG}

% -------------------------------------------------------------------
% INFORMACIÓN DEL TFG Y EL AUTOR
% -------------------------------------------------------------------
\titulo{Inteligencia Artificial Explicable en medicina: métodos para introducir al humano en la adquisición de conocimiento}
\autor{Santiago López Cerro}
\grado{Doble Grado en Ingeniería Informática y Matemáticas}
\facultad{Facultad de Ciencias y Escuela Técnica Superior de Ingenierías Informática y de Telecomunicación}
\universidad{Universidad de Granada}
\curso{2026--2027}
\tutor{
  Francisco Herrera Triguero \\ \emph{Departamento de Ciencias de la Computación e Inteligencia Artificial }

  % Añadir tantos tutores como sea necesario.

  \medskip
  Iván Sevillano García \\ \emph{Departamento de Ciencias de la Computación e Inteligencia Artificial }
}

\begin{document}

\maketitle

% -------------------------------------------------------------------
% FRONTMATTER
% -------------------------------------------------------------------
\frontmatter % Desactiva la numeración de capítulos y usa numeración romana para las páginas

% !TeX root = ../tfg.tex
% !TeX encoding = utf8
%
%*******************************************************
% Declaración de originalidad
%*******************************************************

\thispagestyle{empty}
\makeatletter

\hfill\vfill

\textsc{Declaración de originalidad}\\\bigskip

D./Dña. \@author \\\medskip

Declaro explícitamente que el trabajo presentado como Trabajo de Fin de Grado (TFG), correspondiente al curso académico \ugrTFG@academicYear, es original, entendido esto en el sentido de que no he utilizado para la elaboración del trabajo fuentes sin citarlas debidamente.
\medskip

En Granada a \today
\vspace{3cm}
\begin{center}
  Fdo: \@author

\end{center}

\vfill

\makeatother
\cleardoublepage
\endinput

% !TeX root = ../tfg.tex
% !TeX encoding = utf8

%*******************************************************
% Dedication
%*******************************************************
\thispagestyle{empty}
\phantomsection 
\pdfbookmark[1]{Dedicatoria}{Dedicatoria}

\hfill
\vfill

\begin{flushright}
\itshape
Dedicatoria (opcional) \\
Ver archivo \texttt{preliminares/dedicatoria.tex}
\end{flushright}

\vfill

\cleardoublepage
\endinput
                % Opcional
% !TeX root = ../tfg.tex
% !TeX encoding = utf8

%*******************************************************
% Agradecimientos
%*******************************************************

\chapter{Agradecimientos}

\pendiente{agradecimientos}

\cleardoublepage
\endinput
            % Opcional
% !TeX root = ../tfg.tex
% !TeX encoding = utf8
%
%*******************************************************
% Índice de acrónimos (tabla manual)
%*******************************************************

\chapter{Índice de acrónimos}

\begin{longtable}{@{}lp{0.72\textwidth}@{}}
\toprule
\textbf{Acrónimo} & \textbf{Significado} \\
\midrule
\endhead
\bottomrule
\endlastfoot
APS & \emph{Adaptive Prediction Sets} \\
ASH & \emph{Activation Shaping} \\
AUPR & Área bajo la curva precisión-exhaustividad (\emph{Area Under the Precision-Recall curve}) \\
AUROC & Área bajo la curva ROC (\emph{Area Under the Receiver Operating Characteristic curve}) \\
CCR & Tasa de clasificación correcta (\emph{Correct Classification Rate}) \\
CNN & Red neuronal convolucional (\emph{Convolutional Neural Network}) \\
CRP & \emph{Concept Relevance Propagation} \\
DICE & \emph{Directed Sparsification} \\
DTD & \emph{Describable Textures Dataset} \\
EBO & \emph{Energy-Based OOD} \\
FPR & Tasa de falsos positivos (\emph{False Positive Rate}) \\
H\&E & Hematoxilina y eosina \\
ID & Dentro de la distribución (\emph{In-Distribution}) \\
KNN & $k$ vecinos más cercanos (\emph{$k$-Nearest Neighbors}) \\
LAC & \emph{Least Ambiguous set-valued Classifier} \\
LDA & Análisis discriminante lineal (\emph{Linear Discriminant Analysis}) \\
LOSO & \emph{Leave-One-Slide-Out} \\
LRP & \emph{Layer-wise Relevance Propagation} \\
MDS & \emph{Mahalanobis Distance Score} \\
MLS & \emph{MaxLogit Score} \\
MSP & \emph{Maximum Softmax Probability} \\
OOD & Fuera de la distribución (\emph{Out-of-Distribution}) \\
OSR & Reconocimiento en conjunto abierto (\emph{Open Set Recognition}) \\
PCA & Análisis de componentes principales (\emph{Principal Component Analysis}) \\
RAPS & \emph{Regularized Adaptive Prediction Sets} \\
ReAct & \emph{Rectified Activations} \\
RMDS & \emph{Relative Mahalanobis Distance Score} \\
SHE & \emph{Simplified Hopfield Energy} \\
STOOD-X & \emph{Statistical nonparametric Test for OOD Detection with eXplainability} \\
TPR & Tasa de verdaderos positivos (\emph{True Positive Rate}) \\
ViM & \emph{Virtual-logit Matching} \\
WSI & Imagen de lámina completa (\emph{Whole Slide Image}) \\
XAI & Inteligencia artificial explicable (\emph{eXplainable Artificial Intelligence}) \\
\end{longtable}

\endinput


% !TeX root = ../tfg.tex
% !TeX encoding = utf8
%
%*******************************************************
% Summary (800--1500 palabras, en inglés)
%*******************************************************

\selectlanguage{english}
\chapter{Summary}

% Puntos que debería cubrir (se escribe al final):
% - Context: deep learning for prostate histopathology classification (14 classes, H&E patches).
% - Problem: models deployed in clinics receive inputs outside the training distribution (OOD).
% - Research questions and hypotheses.
% - Mathematical part: formalisation of OOD detection, ROC/AUROC and its link with the
%   Mann-Whitney U statistic, post hoc detectors (logits, activations, features/distances),
%   Mann-Whitney test, bootstrap/McNemar, conformal prediction (marginal coverage theorem,
%   APS/LAC, Mondrian), LRP/CRP.
% - Computer science part: five architectures compared, Swin-S selected; 16 OOD detectors
%   evaluated on near-OOD (NCT-CRC), internal control, covariate shift (PANDA), resolution
%   control and far-OOD (DTD); STOOD-X reimplementation; conformal prediction; slide leakage.
% - Main results and conclusions; limitations; future work (explainability and bias removal).
% - Keywords.
\pendiente{summary en inglés (800--1500 palabras)}

% Al finalizar el resumen en inglés, volvemos a seleccionar el idioma español para el documento
\selectlanguage{spanish}
\endinput

% !TeX root = ../tfg.tex
% !TeX encoding = utf8
%
%*******************************************************
% Resumen en castellano
%*******************************************************

\chapter{Resumen}

% Puntos que debería cubrir:
% - Contexto y problema (clasificación histopatológica de próstata y detección OOD).
% - Objetivos de la parte matemática y de la parte informática.
% - Metodología y principales resultados.
% - Conclusiones y trabajo futuro.
% - Palabras clave.
\pendiente{resumen en castellano}

\endinput


% !TeX root = ../tfg.tex
% !TeX encoding = utf8

%*******************************************************
% Table of Contents
%*******************************************************
\phantomsection
\pdfbookmark[0]{\contentsname}{toc}

\setcounter{tocdepth}{2} % <-- 2 includes up to subsections in the ToC
\setcounter{secnumdepth}{3} % <-- 3 numbers up to subsubsections

\tableofcontents
\label{sec:tableofcontents}

%*******************************************************
% List of Figures and of the Tables
%*******************************************************

% *******************************************************
%  List of Figures
% *******************************************************
\phantomsection
% \listoffigures

%*******************************************************
% List of Tables
%*******************************************************
\phantomsection
% \listoftables

%*******************************************************
% List of Listings
% The package \usepackage{listings} is needed
%*******************************************************
% \phantomsection
% \renewcommand{\lstlistlistingname}{Listados de código}
% \lstlistoflistings

\cleardoublepage


% -------------------------------------------------------------------
% MAINMATTER
% -------------------------------------------------------------------
\mainmatter % activa la numeración de capítulos, resetea la numeración de las páginas y usa números arábigos

% !TeX root = ../tfg.tex
% !TeX encoding = utf8

\chapter{Introducción}\label{cap:introduccion}

\pendiente{introducción}

\endinput


% -------------------------------------------------------------------
% BACKMATTER
% -------------------------------------------------------------------

\backmatter % Desactiva la numeración de los capítulos
\pdfbookmark[-1]{Referencias}{BM-Referencias}

% -------------------------------------------------------------------
% BIBLIOGRAFÍA
%-------------------------------------------------------------------

\bibliographystyle{plain-es}
% Otros estilos bibliográficos: alpha-es, amsalpha, plain
\begin{small}
  \bibliography{library.bib}
  \label{sec:bibliography}
\end{small}

\end{document}
